\documentclass[11pt]{article}
\usepackage[margin=1in]{geometry}
\usepackage{booktabs}
\usepackage{hyperref}

\title{TPCRP Coursework Progress Draft}
\author{Kasim}
\date{\today}

\begin{document}
\maketitle

\section{Project Goal}
The first task of the coursework is to implement the TPCRP active learning method from Hacohen et al.\ for CIFAR-10. The algorithm is a low-budget active learning strategy that selects representative and diverse unlabeled samples for annotation. Our current goal is to reproduce the paper's baseline method faithfully and ensure that it runs end-to-end without implementation errors.

\section{Paper Understanding}
The paper argues that low-budget active learning should prefer typical, representative examples rather than highly uncertain ones. The TPCRP variant used for CIFAR-10 first learns image representations with SimCLR, then clusters the embeddings, and finally queries the most typical samples from the most underrepresented clusters.

The selection rule implemented so far is:
\begin{enumerate}
    \item Train a self-supervised encoder on CIFAR-10.
    \item Extract L2-normalized embeddings for the training pool.
    \item Cluster the embeddings with K-means or MiniBatchKMeans.
    \item Prioritize clusters with the fewest labeled samples.
    \item Break ties by choosing the largest cluster.
    \item Within that cluster, compute typicality as the inverse mean distance to the $k$ nearest neighbours.
    \item Query the unlabeled point with the highest typicality.
\end{enumerate}

\section{Implementation Completed So Far}
The repository now contains a working baseline structure for Task 1:
\begin{itemize}
    \item \texttt{models.py}: a CIFAR-adapted ResNet-18 SimCLR encoder and a supervised ResNet-18 classifier.
    \item \texttt{utils.py}: CIFAR-10 loading, SimCLR augmentations, NT-Xent loss, embedding extraction, and training/evaluation utilities.
    \item \texttt{tpcrp.py}: the TPCRP query selector with clustering and typicality scoring.
    \item \texttt{train.py}: an end-to-end pipeline that trains representations, selects query samples, trains a supervised classifier on the labeled set, and records metrics.
\end{itemize}

The current implementation follows the main TPCRP logic from the paper and uses paper-leaning defaults for SimCLR training on CIFAR-10, including a batch size of 512, learning rate 0.4, weight decay $10^{-4}$, cosine scheduling, and a maximum of 500 clusters.

\section{Current Outputs}
Each run now saves experiment artefacts into an \texttt{outputs/} directory:
\begin{itemize}
    \item \texttt{round\_\#\#\_metrics.csv}: per-epoch supervised training metrics.
    \item \texttt{run\_summary.json}: selected query indices and round-level accuracy summaries.
\end{itemize}
These files will be useful later for the report tables, plots, and reproducibility evidence.

\section{Current Status and Remaining Work}
The baseline algorithm structure is implemented, but full empirical validation is still pending. The next steps are:
\begin{enumerate}
    \item install the dependencies;
    \item run a short smoke test to verify the full training pipeline;
    \item run a longer experiment with more realistic TPCRP settings;
    \item compare accuracy across rounds and use the saved outputs for the report.
\end{enumerate}

\section{Notes for the Final Report}
This draft only documents implementation progress. The final two-page report should later add:
\begin{itemize}
    \item a concise problem introduction;
    \item a clearer methodology section tied directly to the paper;
    \item tables and plots from actual experimental runs;
    \item a discussion of the modified TPCRP variant and its impact.
\end{itemize}

\end{document}
